
Let us recall some notations.
Let $X = [x_0, x_1, \dots x_{n-1}]$ of sequence of $n$ distinct objects and 
$W =[w_0, \dots, w_{n-1}]$ a sequence of non-negative integers such that 
object $x_i$ is of weight $w_i$. We denote the total weight by $w = \sum_i w_i$.

We aim at generating at random an object $x_i$
according to the global distribution $W$ avoiding any bias: the probability for generating $x_i$ is given 
by $w_i / w$, for each $i$.

For the rest of the paper, we suppose 
that the weights are positive integers, otherwise we remove all objects $x_i$
of weight $0$. Furthermore we suppose that the greatest common divisor of integers in $W$ is $\texttt{GCD}(W) = 1$.
If it is not the case, we replace $W$ by the normalized sequence 
$[\frac{w_0}{\texttt{GCD}(W)}, \dots, \frac{w_{n-1}}{\texttt{GCD}(W)}]$,
still being an integer distribution. 
Obviously both sets of weights (the original one and the latter one) yield the same probability distribution. 
So, from now we suppose $W$ contains only positive integers whose $\texttt{GCD}$ is $1$.

\begin{definition}
    We define the cell size, denoted by $\cs$, to be the integer $\lfloor \frac{w}{n} \rfloor$.
\end{definition}

\begin{remark*}
    %The number $\lfloor \frac{w}{n} \rfloor$ is equal to $w // n$.
    Since the distribution $W$ contains $n$ positive integers, we have $\cs \geq 1$.
\end{remark*}

We compute the Euclidean division of $w$ by $\cs$ and thus define the quotient $q$ and the remaining part $r$ such that
\begin{equation}
    \label{eq:division}
    w = q \cdot \cs + r \qquad\qquad \text{with} \quad 0 \leq r < \cs.
\end{equation}
Let us focus on the simple running example where $X$ contains $n=5$ objects from $x_0$ to $x_4$ 
with respective weights $W = [2, 9, 1, 6, 1]$.
The total weight is $w = 19$ and the cell size is $\cs = 3$. Thus we calculate $q = 6$ and $r=1$.\\
From now, generally, we are faced in two cases:
\begin{itemize}
    \item if $r > 0$, then we define a new object $x_n$, called \emph{virtual object}
        and relate it to the weight $w_n = \cs - r$.
        Furthermore, we define $\bar{X} = [x_0, x_1, \dots x_{n-1}, x_n]$, 
        $\bar{W} =[w_0, \dots, w_{n-1}, w_n]$ and $\bar{w} = w + w_n$.
    \item If $r=0$, the we choose $\bar{X} = X$ and $\bar{W} = W$, and we set $x_n=\emptyset$.
\end{itemize}
Thus, we will use the following notations
\begin{equation}
    \label{eq:division2}
    \bar{w} = \bar{q} \cdot \cs \qquad\qquad \text{with either } \bar{q}=q+1 \text{ or } \bar{q}=q.
\end{equation}

\begin{remarks*}
    We observe now that $\bar{w}$ is a multiple of $\cs$.
    The value $\cs$ does not change, even if $\bar{X} \neq X$.
\end{remarks*}
In our running example, we append a virtual object $x_5$ with weight $2$ and then $\bar{q} = 7$.

Next algorithm builds a table $T$ with $\bar{q}$ cells,
each one containing a triplet $(v', x', x)$ where $v'$ is an integer in 
$\llbracket 1, \cs \rrbracket$\footnote{The notation $\llbracket a, b \rrbracket$ 
    stands for the subset of all consecutive integers from $a$ to $b$.}
and $x$ and $x'$ being two distinct objects from $\bar{X}$. The object $x$ may eventually be $\emptyset$
in the case where $v'=\cs$, in that case we will say that $x'$ has no alias.
As in the classical Alias method, 
the preprocessing step objective is to arrange the weight distribution.
Here the weights $\bar{W}$ must be distributed into $\bar{q}$ cells,
with the constraint that each cell contains at most two different objects whose associated weights cumulate to $\cs$.

Intuitively, we apply Vose's preprocessing step with specific inputs that are different from the original ones.
Let us present our algorithm approach for this preprocessing step, building the alias table $T$.
The algorithm proceeds in two phases. 
First objects are distributed in two stacks $L$ and $S$,
containing respectively the larger objects (whose individual weights
are larger than $\cs$)
and the smaller objects of weights smaller than or equal to $\cs$. Then
(1) the first step of the algorithm is to empty the stack $L$.
To that order, each object of $S$ is completed with a part of a larger object from $L$
if necessary. In the case where the smaller object is of weight $\cs$ it is put alone, without alias, in its cell.
Once done,  either $S$ is also empty and the algorithm terminates (correctly)
or from there (2) we will prove that $S$ will contain objects whose remaining weight are exactly $\cs$. Hence each object is put
without alias in a new cell of $T$.
This approach will be proved correct. 

Taking our running example, we get the following result for $T$:
\[
    [(2, x_5, x_3), (1, x_4, x_3), (3, x_3, \emptyset), (1, x_2, x_1), (2, x_0, x_1), (3, x_1, \emptyset), (3, x_1, \emptyset) ].
\]

\begin{algorithm}[htb]
\caption{Building of the table of triplets $T$} 
\label{algo:table}
% Old version
	% \begin{algorithmic}[1]
    %     \renewcommand{\algorithmicrequire}{\textbf{Input:}}
    %     \renewcommand{\algorithmicensure}{\textbf{Output:}}
    %     \Require $\bar{X}$, $\bar{W}$, $\cs, \bar{q}, x_n$ previously defined.
    %     \Ensure  The table of triplets $T$.
    %     \Function{TableBuilding}{$\bar{X}$, $\bar{W}$, $\cs, \bar{q}, x_n$}
    %         \State $T = \texttt{init}(\bar{q})$ \Comment{Initialization of a table indexed from $0$ to $\bar{q}-1$}
    %         \State $L = \{x \in \bar{X} \quad | \quad W[x] \geq \cs\}$ \Comment{$L$ is the stack of the larger objects}
    %         \State $S = \{x \in \bar{X} \quad | \quad W[x] < \cs\}$ \Comment{$S$ is the stack of the smaller objects}
    %         \State $k = 0$
    %         \While{$L$ is not empty}
    %             \State $x = \texttt{pop}(L)$ \Comment{pop operation extracts one object from the stack}
    %             \If{$S$ is not empty}
    %                 \State $x' = \texttt{pop}(S)$
    %                 \If{$x' \neq x_n$} \Comment{$x'$ is not the virtual object}
    %                     \State $T[k] = (\bar{W}[x'], x', x)$
    %                 \Else \Comment{$x'$ is the virtual object}
    %                     \State $T[\bar{q}-1] = (\cs-\bar{W}[x'], x, x')$ \Comment{$x'$ is the 3rd component of the triplet, in the last cell}
    %                     \State $k = k - 1$
    %                 \EndIf
    %                 \State $\bar{W}[x] = \bar{W}[x] - (\cs - \bar{W}[x'])$
    %             \Else
    %                 \State $T[k] = (\cs, x, \emptyset)$
    %                 \State $\bar{W}[x] = \bar{W}[x] - \cs$
    %             \EndIf
    %             \If{$0 < \bar{W}[x] < \cs$}
    %                 \State $push(S, x)$ \Comment{push operation adds one object to the stack}
    %             \ElsIf{$\bar{W}[x] \geq \cs$}
    %                 \State $push(L, x)$
    %             \EndIf
    %             \State $k = k + 1$
    %         \EndWhile
    %         \State \Return $T$
    %     \EndFunction
	% \end{algorithmic}
%New version
	\begin{algorithmic}[1]
        \scriptsize
        \renewcommand{\algorithmicrequire}{\textbf{Input:}}
        \renewcommand{\algorithmicensure}{\textbf{Output:}}
        \Require $\bar{X}$, $\bar{W}$, $\cs, \bar{q}$ previously defined.
        \Ensure  The table of triplets $T$.
        \Function{TableBuilding}{$\bar{X}$, $\bar{W}$, $\cs, \bar{q}$}
            \State $T = \texttt{init}(\bar{q})$ %\Comment{Initialization of a table indexed from $0$ to $\bar{q}-1$}
            \State $L = \{x \in \bar{X} \quad | \quad \bar{W}[x] > \cs\}$ \Comment{$L$ is the stack of the larger objects}
            \State $S = \{x \in \bar{X} \quad | \quad \bar{W}[x] \leq \cs\}$ \Comment{$S$ is the stack of the smaller objects}
            \State $k = 0$
            \While{$L$ is not empty}
                \State $x = \texttt{pop}(L)$ %\Comment{pop operation extracts one object from the stack}
                \If{$S$ is not empty}
                    \State $x' = \texttt{pop}(S)$
                    \If{$\bar{W}[x'] < \cs$}
                        \State $T[k] = (\bar{W}[x'], x', x)$
                        \State $\bar{W}[x] = \bar{W}[x] - (\cs - \bar{W}[x'])$
                    \Else
                        \State $T[k] = (\cs, x', \emptyset)$
                    \EndIf
                \Else
                    \State $T[k] = (\cs, x, \emptyset)$
                    \State $\bar{W}[x] = \bar{W}[x] - \cs$
                \EndIf
                \If{$0 < \bar{W}[x] \leq \cs$}
                    \State $push(S, x)$ %\Comment{push operation adds one object to the stack}
                \ElsIf{$\bar{W}[x] > \cs$}
                    \State $push(L, x)$
                \EndIf
                \State $k = k + 1$
            \EndWhile
            \For{$x' \in S$}
                \State $T[k] = (\cs, x', \emptyset)$ \Comment{the objects still in $S$ have a remaining weight equal to $\cs$}
                \State $k = k + 1$
            \EndFor
            \State \Return $T$
        \EndFunction
	\end{algorithmic}    
 \end{algorithm}

\begin{remark*}
    The conditional cases in Lines 10 to 14 could be omitted by keeping only Line 12 with $W[x']=\cs$; thus 
    the alias object $x$ would have a weight 0 in this cell. The algorithm would still be correct.
\end{remark*}


\begin{proposition}
\label{prop:size_table}
    The size $\bar{q}$ of the table $T$ satisfies $n \leq \bar{q} \leq (1 + \frac{1}{\cs})\cdot n -1$.
\end{proposition}
\begin{proof}
    First suppose $q = \bar{q}$ and $w = \bar{w}$, i.e. $r=0$ in Equation~\eqref{eq:division2}. In that case,
    using the latter equation, and the definition of $\cs$ we have
    \[
        \frac{w}{q} = \lfloor \frac{w}{n} \rfloor.
    \]
    Thus we deduce $q \geq n$.
    Furthermore, suppose now that $q \geq (1 + \frac{1}{\cs})\cdot n$.
    \[
        \cs = \lfloor \frac{q\cdot \cs}{n} \rfloor \geq \lfloor \frac{(1 + \frac{1}{\cs})\cdot n \cdot \cs}{n} \rfloor \geq \cs +1.
    \]
    This is a contradiction.
    So in this case $n \leq q < (1 + \frac{1}{\cs})\cdot n$.\\
    Second, suppose in the Equation~\eqref{eq:division}, we are in the case $r>0$. From Equation~\eqref{eq:division2}
    we obtain
    \[
        \frac{w + (\cs - r)}{q+1} = \lfloor \frac{w}{n} \rfloor.
    \]
    Thus we get
    \[
        \frac{w}{q+1} \leq \lfloor \frac{w}{n} \rfloor \leq \frac{w}{n},
    \]   
    so $q+1\geq n$.
    Again, suppose $q \geq (1 + \frac{1}{\cs})\cdot n -1$. We deduce
    \[
        \cs = \lfloor \frac{(q+1)\cdot \cs}{n} \rfloor \geq \lfloor \frac{ (1 + \frac{1}{\cs})\cdot n \cdot \cs}{n} \rfloor 
            \geq \lfloor (1 + \frac{1}{\cs}) \cs \rfloor = \cs + 1.
    \]    
    This is a contradiction so we obtain in this second setting $n-1 \leq q < (1 + \frac{1}{\cs})\cdot n-1$.\\
    Using the relationship between $q$ and $\bar{q}$ from~\eqref{eq:division2}, we get the proposition's statement.\qed
\end{proof}
Note that these bounds for $\bar{q}$ are tight. In fact, if all the weights are equal (thus we are faced to the uniform distribution)
once normalized each weight is 1, then $q=\bar{q} = n$.
For the upper bound, fix two positive integers $\omega$ and $n$, and take a set of $\omega \cdot n$ objects, such that
all of them except one have weight $\omega$ and the last one has weight $\omega\cdot n$.
We calculate that the constant $\cs$ is equal to $\omega$, thus we get
$q = (1 + \frac{1}{\cs})\cdot \omega\cdot n - 1$ (recall $\omega\cdot n$ is the number of objects).


\begin{lemma}
\label{lem:seq_increase}
    After each loop, starting Line 6 of Algorithm~\ref{algo:table}, the average updated weight of the objects still lying either in $L$
    or in $S$ is greater than or equal to $\cs$.
\end{lemma}
\begin{proof}
    We prove the result by induction.
    \begin{itemize}
        \item At the beginning of the algorithm, let us prove that the average weight of the objects is larger than or equal to $\cs$.
            \begin{itemize}
                \item In the case when $r=0$ in Equation~\eqref{eq:division},
                    we have $\bar{w} = w$ and the number of objects is $n$, thus the average weight is $\bar{w}/n = w/n \geq \cs$.
                \item In the case when $r>0$, we have $\bar{w} = w + (\cs - r) = (q+1)\cs$ and the number of objects is $n+1$, thus the average weight satisfies
                    \[
                        \frac{w}{n+1} < \frac{\bar{w}}{n+1} = \frac{(q+1)\cs}{n+1} < \frac{q+1}{n+1}\cdot\frac{w}{n} = \frac{q+1}{n}\cdot \frac{w}{n+1}.
                    \]
                    \Cref{prop:size_table} induces $n\leq q+1$, using the later inequality the case $n=q+1$ is not possible,
                    thus
                    \[
                         \frac{\bar{w}}{n+1} = \frac{(q+1)\cs}{n+1} \geq \frac{(n+1)\cs}{n+1} = \cs.
                    \]
            \end{itemize}
            Hence when we enter the loop Line 6, the average weight of the objects distributed in $L$ and $S$ is greater than or equal to $\cs$.
        \item Suppose at the beginning of the loop, it remains $k$ objects distributed in $L$ and $S$, 
            whose average weight $\omega \geq \cs$.\\
            During the loop, the total weight is decreased by $\cs$.\\
            There are two mutually exclusive loop-behaviors:
            \begin{itemize}
                \item An object from $S$ is removed from $S$;
                \item An object from $L$ is of updated weight greater than $\cs$, its new weight is decreased by $\cs$
                    and then, either it still stands in $L$ or it is pushed in $S$.
            \end{itemize}
            In the first case, at the end of the loop, the cumulative remaining weight is $W' = k\cdot \omega - \cs$,
            for $k-1$ objects (since $\omega$ was the average weight of the $k$ objects). 
            If the number of remaining objects $k-1$ is positive, then the average weight is
            \[
                \omega' = \frac{W'}{k-1} = \frac{k\cdot \omega - \cs}{k-1}  \geq \frac{k\cdot \cs - \cs}{k-1} = \cs.
            \]
            Otherwise, when $k-1 = 0$, the algorithm terminates.\\

            In the second case, the following constraints are valid : $S$ is empty and there is an element in $L$
            of weight greater than $\cs$.
            Thus the total weight $\bar{W}$ of objects remaining in $L$ at the beginning of the loop is at least $k\cdot (\cs + 1) = k\cs + k$.
            But the remaining weight is a multiple of $\cs$, thus $\bar{W} \geq (k+1) \cdot \cs$,
            since the number $k$ of objects does not decrease, so we deduce the updated average weight
            \[
                \frac{W'}{k} = \frac{\bar{W} - \cs}{k} \geq \frac{(k+1)\cdot \cs - \cs}{k}  = \cs.
            \]
            So in both cases, at the end of the loop, the average weight of the remaining objects is at least $\cs$.
    \end{itemize}
    Thus the lemma is proved.\qed
\end{proof}
\begin{lemma}
\label{lem:res_without aliases}
    At Line 27 of Algorithm~\ref{algo:table}, all objects still lying in $S$ have a weight equal to $\cs$.
\end{lemma}
\begin{proof}
    Arriving at Line 27, the loop invariant proved in~\Cref{lem:seq_increase} is still valid, thus the 
    average weight of the objects still lying in $S$ is greater than or equal to $\cs$.
    Hence we deduce that these remaining objects have all a weight equal to $\cs$.\qed
\end{proof}


\begin{theorem}
\label{theo:correctness}
    Algorithm~\ref{algo:table} is correct.
\end{theorem}
The correctness of the approach lies essentially in proving that the algorithm terminates, $T$ being completed.
\begin{proof}
    This is a direct corollary of~\Cref{lem:seq_increase,lem:res_without aliases}. Each cell of $T$ contains
    a cumulated weight of $\cs$, either using a single object or by completing an object with an alias.\qed
\end{proof}



% \noteAnt{voir si on garde la suite, je pense que ce qui est indiqué n'est pas clair}

% In the next proposition, we propose a deeper analysis of the size of the alias table.
% Beyond Equation~\eqref{eq:division}, let us introduce a second division:
% \begin{equation}
%     \label{eq:division3}
%     w = \cs \cdot n + \tilde{r} \qquad\qquad \text{with} \quad 0 \leq \tilde{r} < n.
% \end{equation}
% Note, in the division of $w$ by $n$, the quotient is $\cs$.
% \begin{proposition}
%     The number of supplementary cells $q-n$ satisfies $(q-n)\cs = \tilde{r} - r$,
%     with values $r$ and $\tilde{r}$ respectively given by Equations~\eqref{eq:division} and~\eqref{eq:division3}.
% \end{proposition}
% \begin{proof}
%     Using Equations~\eqref{eq:division} and~\eqref{eq:division3}, we get
%     two constants $0\leq r < \cs$ and $0 \leq \tilde{r} < n$.
%     We get the two equalities $w + \cs - r = (q+1)\cdot\cs$ and $w = n \cdot \cs + \tilde{r}$.
%     We easily deduce $(q-n)\cdot\cs = \tilde{r} - r$.
% \end{proof}
% \noteAnt{intuitivement, avec un peu d'arithmétique, je pense qu'on peut trouver toutes les combinaisons
% satisfaisant cette contrainte de divisibilité.}